\section{Larger fixed isotropic blocks}\label{app:fixedL}
Fix an odd $L\geq5$ before letting $K$ grow. In a block of length $L$,
\eqref{eq:Blocal} contains $2L-2$ factors, so the quadratic coefficients
are bounded by $H=K(4K)^{2L-2}$. The alternating diagonal form has
balanced signature. We need a bound for a whole isotropic collection,
rather than repeated applications of a bound for one vector without
tracking the intervening basis changes.

For a nonsingular integral symmetric $N$-by-$N$ matrix of balanced
signature, odd $N\geq5$, with entries bounded by $H\geq1$, set
\[
 V_N=(3N^2H)^{(N-1)/2},\quad F_N=8N^2H^2V_N^2,
 \quad H'_N=N^2HF_N^2.
\]
Define $T_5(H)=V_5$ and
$T_N(H)=\max(V_N,NF_NT_{N-2}(H'_N))$. For fixed $N$ this is
polynomial in $H$. We claim that $(N-3)/2$ independent mutually
orthogonal integral isotropic vectors can be chosen with sup norm at
most $T_N(H)$.

For $N=5$ this is Meyer and Cassels. In higher dimension choose a
nonzero isotropic $v$ with norm at most $V_N$, and a coordinate $j$
such that $b=v^TAe_j\ne0$. Put $w=2be_j-A_{jj}v$. Then
$w^TAw=0$ and $v^TAw=2b^2$, so $v,w$ split a rational hyperbolic
plane. The integral vectors
\[
 f_i=2b^2e_i-(e_i^TAw)v-(e_i^TAv)w
\]
span its orthogonal complement and have coordinates bounded by
$(6N^2+2N)H^2V_N^2\leq F_N$. Select $N-2$ independent columns to
form $F$. The induced matrix $F^TAF$ is nonsingular, balanced, and
bounded by $H'_N$. Apply induction there and lift back by $F$,
retaining $v$ as the first isotropic vector. This proves the claimed
recursion and controls the entire collection.

Applying this construction in the fixed $L$-blocks and then using
\eqref{eq:full-lift} gives integral cofactor coordinates of size
$K^{O_L(1)}$. Hadamard bounds $\det(\mathcal U^T\mathcal U)$ even
when supports overlap within a block. Cauchy--Binet and
$\sum_I|\det\mathcal U_I|\leq
\sqrt{\binom{2K}{h}\det(\mathcal U^T\mathcal U)}$
give $O_L(K\log K)$ logarithmic transfer cost. The resulting dimension
is $K+(L-3)\lfloor K/L\rfloor/2$. No uniformity for growing $L$
is asserted.

\section{The normalized-decay ceiling}\label{app:ceiling}
Let $P$ be any real polynomial of degree below $K$, with
\[
 G(P)=\int_0^\infty\frac{P(4y^2)^2w(y)}{D(4y^2)}\dd y,
 \qquad
 B(P,P)=\sum_{l\ \mathrm{odd}}\frac{b_lP(-l^2)^2}{D'(-l^2)}.
\]
The weight bound in Lemma~\ref{lem:gramlower} gives
$w(y)\geq c_0e^{-\pi y}$ with $c_0=1/(48\pi)$.
Put $p(t)=P(K^2t)$. Restriction to $0\leq y\leq K/2$ yields
\[
 G(P)\geq\frac{c_0K e^{-\pi K/2}}{2D(K^2)}
                         \int_0^1p(x^2)^2\dd x.
\]
Let $\mathsf P_j$ be the standard Legendre polynomial, normalized by
$\mathsf P_j(1)=1$. The shifted functions
$\phi_j(t)=\sqrt{2j+1}\mathsf P_j(2t-1)$ are orthonormal for Lebesgue
measure on $[0,1]$. The classical integral formula at $x>1$
\cite[14.12.7]{DLMF}, together with parity and continuity, gives
$|\mathsf P_j(-1-2a^2)|\leq(9+\sqrt{80})^j$ for $0\leq a\leq2$.
Also $\int_0^1p(t)^2\dd t\leq2\int_0^1p(x^2)^2\dd x$.
Since $\sum_{j=0}^{K-1}(2j+1)=K^2$, Cauchy--Schwarz in this
orthonormal basis gives
\[
 |P(-l^2)|^2\leq\frac{4K}{c_0}e^{\pi K/2}D(K^2)
       (2+\sqrt5)^{4K-4}G(P).
\]
Let $N=2K$. Formula \eqref{eq:Dprime} implies
$\min_l|D'(-l^2)|\geq(N!)^2/(2N^2)$. The increasing function
$\log(1/4+x^2)$ has right-sum error at most $\log5$. Since its
integral on $[0,1]$ is $\log(5/4)-2+\arctan2$ and
$N!\geq(N/e)^N$,
\[
 \frac{D(K^2)}{\min_l|D'(-l^2)|}
 \leq40K^2e^{2K(\log(5/4)+\arctan2)}.
\]
Finally $\sum_{l\ \mathrm{odd}}|b_l|\leq4K^4$. Thus
\[
 |B(P,P)|\leq30720\pi K^7e^{C_{1/2}K}G(P),\quad
 C_{1/2}=\pi/2+2\log(5/4)+2\arctan2+4\log(2+\sqrt5).
\]

For any compressed positive matrix $G$ and beta coefficient $B_U$,
the symmetric operator $G^{-1/2}B_UG^{-1/2}$ has eigenvalues bounded
in absolute value by this same quantity. If its nonzero eigenvalues
are $\lambda_1,\ldots,\lambda_h$, then
\[
 D_U(X)=D_U(\bet)\prod_{i=1}^h(1+(X-\bet)\lambda_i),\qquad
 \delta_U=\prod_{i=1}^h|\lambda_i|^{-1}.
\]
This proves $\deg D_U=\rank B_U=h$, rather than assuming full rank.
For $h=0$ the normalized value is one. For $h\leq K/5$ the operator
bound gives \eqref{eq:ceiling}. Its strict separation from
\eqref{eq:arith-rate} can be checked without decimals:
$C_{1/2}/5<289/140$ using $\pi<22/7$, $\log(5/4)<1/4$,
$\arctan2<9/8$, $\log(2+\sqrt5)<3/2$, while
$(197+192\log2)/150>13/6$ using $\log2>2/3$.

\section{Evaluation in repeated-exponent spaces}\label{app:evaluation}
This is the repeated-exponent form of M\"untz--Laplace orthogonalization;
see \cite[Section~2]{BEZ} for the underlying orthogonal-system method.
We give the needed norm estimate explicitly.
For $\lambda>0$, define
\[
 T_\lambda f(u)=f(u)-2\lambda\int_0^ue^{-\lambda(u-v)}f(v)\dd v.
\]
On decaying exponential polynomials it is an $L^2(0,\infty)$ isometry.
Indeed, with $h=e^{-\lambda u}*f$, the two functions are
$h'+\lambda h$ and $h'-\lambda h$; their squared-norm difference is
$-2\lambda[|h|^2]_0^\infty=0$. Its range is orthogonal to
$e^{-\lambda u}$ by direct integration.

For positive rates $\lambda_1,\ldots,\lambda_N$, allowing repetitions,
the functions
\[
 e_k=T_{\lambda_1}\cdots T_{\lambda_{k-1}}
                         (\sqrt{2\lambda_k}e^{-\lambda_ku})
\]
are orthonormal. Their Laplace transforms are
$\sqrt{2\lambda_k}(s+\lambda_k)^{-1}
\prod_{i<k}(s-\lambda_i)/(s+\lambda_i)$, so they belong to the
exponential-polynomial space with these rates and multiplicities.
Dimension then makes them a basis. Put $\Lambda=\max\lambda_i$
and $\sigma=\sum\lambda_i$. Expanding the finite convolution product
at a complex point $w$, each term with $m$ convolutions costs at most
$|w|^m e^{\Lambda\max(0,-\Re w)}/m!$. Its scalar coefficients sum
to at most $(2\sigma)^m/m!$. Since
$\sum_{m\geq0}a^m/(m!)^2\leq e^{2\sqrt a}$,
Cauchy--Schwarz gives
\begin{equation}\label{eq:complexeval}
 |H(w)|^2\leq2\sigma
 e^{2\Lambda\max(0,-\Re w)+4\sqrt{2\sigma|w|}}\|H\|_2^2.
\end{equation}

For Section~\ref{sec:gram}, take $H(u)=e^{-u/2}F(1+u)$.
There are thirteen copies of each rate $2j+1/2$, $0\leq j<r$.
Thus $\Lambda<2r$ and $\sigma=13r^2-13r/2\leq13r^2$.
Restriction of the original norm to $y\geq1$ gives
$\|H\|_2^2\leq2e\|F\|_\mu^2$.
At $w_l=z_l-1$ the return to $F(z_l)$ contributes the squared
factor $e^{-1}$. In \eqref{eq:complexeval}, $\Re w_l=-1$ and
$|w_l|\leq M_0$. Consequently
$|F(z_l)|^2\leq52r^2e^{(4+4\sqrt{26M_0})r}\|F\|_\mu^2$,
which is the estimate used in Theorem~\ref{thm:residual}.

\section{Quantitative inputs for the eleven linear forms}\label{app:linear-inputs}
This appendix proves Lemma~\ref{lem:linear-inputs}. We distinguish the
coefficient-column clearer used for lattice height from the full
clearer used for the resulting linear forms.

\subsection{Local counts and the arithmetic envelope}
At the pole $t=-a-1/2$, let $c_{i,a}$ be the principal-part coefficients
of $P(2t+n+1)R_n(t)$. For odd $p$ with $p^2>7n$, define
\[
 F_p(a)=\sum_{j=0}^{28}\#\{v:d_j\leq v\leq n-d_j,\ v\equiv a\pmod p\}
 -2\#\{\ell:-2n-1\leq\ell\leq n,\ p\mid2a+1+2\ell\}.
\]
Taylor coefficient extraction gives $v_p(c_{i,a})\geq i-F_p(a)$:
each nonunit numerator factor supplies a power of $p$, every nonunit
denominator factor costs one, and each extracted Taylor degree can
cost at most one more. All nonzero differences have size below $7n<p^2$.
The integer multiplier has integral local coefficients and cannot
worsen the estimate. Isolated pole classes have integral principal
parts without any extraction loss.

If $M_p$ is the maximum of $F_p$ over repeated classes,
$\max(0,M_p+1)$ clears the complete form. For the rational constant,
group the finite tails at the negative integers $-\ell$. The identity
$(PR_n)'(-\ell)=0$ permits either side of $-\ell$ to represent the
same group. If the nonunit class is isolated choose the other side;
otherwise a differentiated evaluation costs $i+1$ powers, leaving
valuation at least $-F_p-1$. In particular no odd prime $p>n$ is
needed. For odd $p^2\leq7n$, put
$\lambda_p=\lfloor\log_p(7n)\rfloor$. At level $p^k$, the excess
of denominator factors over numerator factors is at most
$(16n+25)/p^k+31$: the twenty-nine pole blocks contribute $22n+29$
terms, and the double zero interval removes $6n+4$.
Consequently the constant local coefficient costs at most
\[
 \sum_{k=1}^{\lambda_p}\left(\frac{16n+25}{p^k}+31\right)
 \leq16v_p(n!)+72\lambda_p.
\]
Taylor extraction adds at most $28\lambda_p$, and the finite-tail
evaluation adds at most $30\lambda_p$. The sufficient exponent is
\[
 16v_p(n!)+130\lfloor\log_p(7n)\rfloor.
\]
These excesses sum to $o(n)$. At two, the squared
half-integer numerator factors cost $6n+4$, the other pole factors
at most $22n+29+29\lfloor\log_2n\rfloor$, and extraction at most
$28\lfloor\log_2n\rfloor$. Therefore
$16v_2(n!)+12n+80(\lfloor\log_2n\rfloor+1)$ suffices.

Put $\delta_j=j/116$. For $x\geq1$, use $0\leq y\leq x-1$
when $x<2$, and $0\leq y\leq1$ otherwise. Define
\begin{align*}
 N_j(x,y)&=\#\{h\in\Z:\delta_jx\leq y+h\leq(1-\delta_j)x\},\\
 C(x,y)&=\lfloor y+x+1/2\rfloor-\lfloor y-2x+1/2\rfloor,\\
 M(x)&=\max_y\left(\sum_jN_j(x,y)-2C(x,y)\right),\\
 e(x)&=\max(0,M(x)+1)-16\lfloor x\rfloor.
\end{align*}
This majorizes the excess prime exponent at $x=n/p$.
The elementary counts give $e(x)\leq32+16\{x\}$.
For $x\geq2$, increasing $x$ by $116$ increases the pole count by
$2552$ and twice the zero count by $696$. Since $M(x)\geq16x-31\geq1$,
the positive part causes no change and $e(x+116)=e(x)$.

The ordering, modulo one, of the finitely many affine endpoints
$\delta_jx,(1-\delta_j)x,-x-1/2,2x-1/2$ changes only at rational
pairwise crossings or integer crossings. On each intervening open
interval the envelope is constant. The supplied certificate enumerates
all such cells: $125$ in $[1,2]$ and $12966$ in $[2,118]$.
Summing thirty periods with upward rational rounding and bounding the
tail after $3482$ by $32/3482+8/3481$ gives
\begin{equation}\label{eq:arithA}
 \int_1^\infty e(x)x^{-2}\dd x\leq
 A=\frac{36699777236396068587597012755199}
         {3663299712317355278800988930048}.
\end{equation}
The checks enumerate the defining endpoints and maxima, not just
selected numerical values. For the asymptotic prime sum, apply the
prime number theorem first on a fixed partition. Event points cost
$O(\log n)$; the remote tail is controlled by the displayed majorant
and Chebyshev's bound before the cutoff tends to infinity.

\subsection{Cancellation improves the full clearer}
For a repeated pole class, choose its least address $0\leq a<p$.
Put $x=n/p$, $r=a/p$ and $k=\lfloor x-r\rfloor$, so its addresses
are $a,a+p,\ldots,a+kp$. When $x<2$, repetition requires $r\leq x-1$.
Define the fixed local rational function
\[
 H_{x,r}(u)=\frac{\prod_{h\in\Z+1/2,\ r-2x\leq h\leq r+x}(u-h)^2}
 {\prod_{l=0}^k(u+l)^{b_l}},\qquad
 b_l=\#\{j:\delta_jx\leq r+l\leq(1-\delta_j)x\},
\]
and its net order $F=\sum b_l-2\#\{h\}$. After scaling, the actual
half-integer numerator roots lie in
\[
 [r-2x-1/(2p),\,r+x+1/(2p)]\cap(\Z+1/2).
\]
This contains the ideal interval used in $H_{x,r}$. Each additional
root supplies $p^2$, so it can only remove a surviving local term
when the intended improvement uses at most two powers of $p$.
For a proper rational function $T$ with these poles and principal
parts $T_l$, put
\[
 \mathcal B(T)=\sum_{l=0}^k(-1)^l\operatorname{res}_{-l}T,
 \qquad
 \mathcal Q(T)=\sum_{v=0}^{k-1}(-1)^v\sum_{l=v+1}^kT_l'(-v-1/2).
\]
On a cell where $M=\max F$, suppose $g\in\{1,2\}$, $M\geq g+1$,
and one rational $\kappa$ satisfies
\begin{equation}\label{eq:localcriterion}
 \mathcal Q(u^hH_{x,r})=\kappa\mathcal B(u^hH_{x,r})
 \quad\text{for all }F\geq M-g+1,\quad 0\leq h<g-(M-F).
\end{equation}
Then $g$ powers of $p$ can be removed from the old clearing exponent
for every beta-two-canceling multiplier, apart from the fixed primes
in the denominator of $\kappa$.

Here is the local-to-global justification. Substitute
$t=-a-1/2+pu$ and multiply by $p^M$. Unit-factor inverses have
Taylor coefficients divisible by the appropriate powers of $p$.
Extra numerator zeros cost $p^2$. Modulo $p^g$, the surviving terms
are therefore
\[
 \sum_{0\leq h<g-(M-F)}c_{a,h}p^{M-F+h}u^hH_{x,r}(u),
 \qquad c_{a,h}\in\Z_{(p)}.
\]
For bounded $x$, all local pole differences and half-integer evaluation
distances are units when $p$ is large. Principal-part extraction
commutes with this reduction. The beta-two coefficient is $-4$ times
the sum of the signed local $\mathcal B$ values, scaled by $p^{1-M}$.
The rational constant has the exact grouped formula
\[
 q(P)=-\sum_{\ell=1}^n(-1)^\ell\sum_{a\geq\ell}f_a'(-\ell),
 \qquad f=PR_n,
\]
and either side of $-\ell$ may be used. The nonunit evaluation between
addresses $a+vp,a+(v+1)p$ is at
$\ell=a+(p+1)/2+vp$, corresponding to $u=-v-1/2$.
It gives, up to a common sign, the signed sum of local $\mathcal Q$
values scaled by $p^{-M-1}$. Unit-denominator terms have at least two
powers after this scaling. Thus \eqref{eq:localcriterion} and exact
beta-two cancellation give $g$ powers in the cleared rational constant.
The remaining beta coefficients, of indices at least four, already
have enough powers. This proves the claimed improvement for the same
forms, including their rational constants.

For example, in $n/2<p<15n/29$ the maximal local function is
\[
 H(u)=\frac{\prod_{h=-5/2,-3/2,-1/2,1/2,3/2}(u-h)^2}
                    {u^{29}(u+1)^{29}}.
\]
Reflection makes both functionals vanish on $H(u)(u+1/2)$.
Exact evaluation gives $\mathcal Q(H)/\mathcal B(H)
=-15010513046/139778381$, so $g=2$ works outside its fixed denominator
primes. The broader certificate checks \eqref{eq:localcriterion}
on every relevant pattern in $582$ open cells of $[1,6]$ and obtains
\begin{equation}\label{eq:savingS}
 S=\int_1^6g(x)x^{-2}\dd x=\frac{62795}{199056}.
\end{equation}
Event points and fixed exceptional primes receive no saving. It follows
that a common full clearer satisfies
\begin{equation}\label{eq:fullclearlinear}
 \log D_n^*\leq16\log(n!)+(A+12\log2-S+o(1))n.
\end{equation}

\subsection{Uniform analytic decay}
Let $H(v)=v\log v$ for $v>0$, with $H(0)=0$, and for $u\geq1$ put
\[
 \Phi(u)=2[H(u+2)-H(u-1)]
       -\sum_{j=0}^{28}[H(u+1-\delta_j)-H(u+\delta_j)].
\]
The exact gamma product for $R_n$ gives, uniformly for $t\geq n+1$,
\[
 \log R_n(t)=-16\log(n!)+n\Phi(t/n)+O(\log(n+t)).
\]
Uniformity includes $t=n+1$: the smallest gamma argument is one,
and $\log\Gamma(v)=v\log v-v+O(\log(v+2))$ is uniform for $v\geq1$.
The finite reciprocal-factor sum bounds the logarithmic derivative
by $O(1+\log(n+t))$. Moreover the multiplier and its derivative
are bounded by $C h_n(P)(1+t/n)^{30}$ and
$C h_n(P)n^{-1}(1+t/n)^{29}$.

The phase has a unique maximum near $1.00165$. One verifies
$\Phi''<0$ on $(1,11/10]$ by the rational upper bound
\[
 2/3-20+\sum_j\bigl((1+\delta_j)^{-1}-(21/10-\delta_j)^{-1}\bigr)<0.
\]
For $u\geq11/10$, $\Phi'(u)\leq
2\log((u+2)/(u-1))-22/(u+1)<0$: the negative of this bound has
no interior minimum, is positive at $11/10$, and tends to zero.
For $u\geq4$, the stronger $u\Phi'(u)<-9$ controls the infinite
tail. Rational log enclosures and a bracket for the critical point
give the supplied upper bound
\[
 \Phi_*\leq
 -\frac{4093415547496757252324770508326269837463}
        {170141183460469231731687303715884105728}
 =-24.0589342582\ldots.
\]
The checker independently encloses the maximum with a tangent on the
concave interval. All summands through $t=n$ vanish; the tail bound
absorbs the multiplier's fixed degree. Thus
\[
 |\mathcal J_n(P)|\leq h_n(P)(n!)^{-16}
                       e^{(\Phi_*+o(1))n}.
\]
Combining with \eqref{eq:fullclearlinear} proves the second part of
Lemma~\ref{lem:linear-inputs}, with
$C_*=-A-12\log2-\Phi_*+S>6.0384$.

\subsection{The cheaper column clearer}
The beta-two coefficient is $-4\sum_a(-1)^ac_{1,a}$. At odd large
primes it needs only $\max(0,M_p-1)$ instead of $\max(0,M_p+1)$.
The envelope saving equals two for $x\geq2$. On $[1,2]$ the exact
partition gives integral $107/116$, so the total saving is $223/116$.
At two, put $z=t+1/2$. Multiplication by $2^{6n+4}$ makes the
numerator integral. The even-address and odd-address denominator
products reduce to powers of $z$ and $z+1$. Their resultant is odd,
so partial fractions by these two groups have numerators in
$\Z_{(2)}[z]$. Each group's residue sum is the coefficient of $z^{-1}$
at infinity and is integral at two. The total residue sum is zero
by proper decay, so the alternating sum is twice one group sum.
Consequently the beta-two coefficients have valuation at least
$-6n-1$, uniformly in the multiplier.

A column clearer $E_n$ therefore satisfies
\[
 \log E_n\leq16\log(n!)+(A-223/116-10\log2+o(1))n.
\]
The coefficient estimate uses a different profile from the positive
tail phase $\Phi$. Put $\mathcal H(v)=v\log|v|$ for $v\ne0$,
$\mathcal H(0)=0$, and define, for $0\leq x\leq1$,
\[
 \Psi(x)=2[\mathcal H(1+x)+\mathcal H(2-x)]
 -\sum_{j=0}^{28}
       [\mathcal H(1-\delta_j-x)+\mathcal H(x-\delta_j)].
\]
On the radius-$1/4$ circle about the pole $-a-1/2$, split each product
at its closest integer. Summing changes $\log(1+O(1/|k|))$ costs
$O(\log n)$ per block. Stirling's formula therefore gives, uniformly
in the address $a$ including endpoints,
\[
 \log|R_n(t)|=-16\log(n!)+n\Psi(a/n)+O(\log n).
\]
The profile is symmetric about $1/2$. Away from its finitely many
cusp points,
\[
 \Psi'(x)=2\log\frac{1+x}{2-x}
 +\sum_{j=0}^{28}\log\frac{|1-\delta_j-x|}{|x-\delta_j|}.
\]
On $0<x<1/2$, every $j>0$ term is positive, and the $j=0$ term
dominates the negative numerator contribution because
$(1-x)(1+x)^2-x(2-x)^2=(1-2x)(x^2-x+1)>0$.
Continuity across the cusps gives the midpoint maximum
\[
 \Psi_*=\Psi(1/2)=6\log(3/2)-2\sum_{j=0}^{28}
                  (1/2-\delta_j)\log(1/2-\delta_j).
\]
Cauchy's residue formula then bounds the coefficient column by
$(n!)^{-16}e^{(\Psi_*+o(1))n}$. Clearing and removing its integer gcd
can only improve this estimate. Thus
\[
 H_*=A-223/116-10\log2+\Psi_*
 =24.5217727858\ldots<24.522,
\]
proving the first part of Lemma~\ref{lem:linear-inputs}.

\section{Record coverage and reproducibility}\label{app:records}
The accompanying guide maps the manuscript to its supporting records
and identifies which checks are recomputations, certificate replays,
or retained finite observations. The mathematical claims are separated
as follows.
\begin{center}\small
\begin{tabular}{p{0.29\linewidth}p{0.61\linewidth}}
\toprule
Construction&Conclusion and remaining condition\\\midrule
Monic zeta substitution&Exact all-subset valuation; the globally cleared
$q^{-2}$ certificate is excluded.\\
Full-radical compression&Integral period elimination with $o(K^2)$
relative-height cost; compatible determinant arithmetic and decay open.\\
Low-degree compression&Uniform prime and dyadic content bounds; true
primitive cost versus achieved decay unresolved.\\
Thirteen-value Gram system&Exact rational-block removal and linear-scale
residual decay; primitive arithmetic unresolved.\\
Fourteen-value linear forms&Formal rank, classified identities, and
eleven uniformly small independent rows; collective nonvanishing open.\\
Fixed-direction transport&Cannot shrink a complete collection while
preserving each individual rational coefficient line.\\
Loewner specialization&Nonvanishing hypotheses checked; one compatible
coarse cost comparison inconclusive.\\\bottomrule
\end{tabular}
\end{center}
The final content bounds supersede the weaker intermediate allowances;
the latter are not separate factors to be multiplied into the result.
Finite tables are retained to document the evidence that motivated
the general questions, not as substitutes for their missing limits.
